\documentclass[12pt,a4paper]{article}

\setlength{\topmargin}{0.0in}
\setlength{\oddsidemargin}{0.33in}
\setlength{\textheight}{9.0in}
\setlength{\textwidth}{6.0in}
\renewcommand{\baselinestretch}{1.25}



\usepackage{hyperref}
% \usepackage[utf8]{inputenc}
% \usepackage[T1]{fontenc}
\usepackage{float}
\usepackage{enumitem}
\usepackage{amsmath}
\usepackage{amssymb}
\usepackage{tikz}
\usetikzlibrary{decorations.pathreplacing}
\usepackage{amsfonts}
\usepackage{graphicx}
\usepackage{subcaption}
\usepackage{caption}
\newcommand{\dv}[2]{\frac{\mathrm{d} #1}{\mathrm{d} #2}}
\usepackage{caption}
\captionsetup[figure]{font=small} 
% \captionsetup[table]{font=small}
\usepackage[capitalize,nameinlink]{cleveref}
\usepackage{mathtools}
\newcommand{\defeq}{\vcentcolon=}
\newcommand{\eqdef}{=\vcentcolon}

% Define environments
\newtheorem{theorem}{Theorem}
\newtheorem{definition}{Definition}
\newtheorem{lemma}{Lemma}
\newtheorem{corollary}{Corollary}



\title{Logs from May $25^{\text{th}}$ to June $2^{\text{nd}}$}
\author{}
\date{}
% \author{Maggie McCarthy}
% \date{May 2026}

\begin{document}

\maketitle

\section{Meeting May 25}

\subsection{Summary of last week's work:}




\begin{enumerate}
    \item Last week I finished up a draft of some background notes. This included filling in some omitted proofs and making some notes on pseudospectra.
    \item I then moved onto the Boffi text and worked through the examples in Sections two and three (up to mixed Laplace problem).
\end{enumerate}



\subsection{Meeting Notes}

\begin{itemize}
    \item Understand D.O.F vs. cells. 

    \item When Boffi says $N = ,$ they are setting $maxh = 1/N$ in netgen.

    \item Netgen works better with firedrake and is what is used for adaptivity.

    \item Why try all the different mesh-types in Boffi? This is in the Boffi text. I forgot it in the meeting. \textbf{This is a good reminder to me that I should think more big picture and underdstand the `why' of each example, rather than just trying to perfectly emulate results. } 

    \item Instead of keeping notes in a notebook it is better to LaTeX it up so that it is easier to transition the work to the thesis later on.

    \item This week I need to work on the presentation. The purpose of the presentation is to force people to have started their project, for the examiners to be assigned to each project, and for the examiners to give feedback on the direction of the project. The questions/feedback they give on these presentations will show up in the viva!


    \item The presentation should be structured as follows:
    \begin{enumerate}
        \item Discuss the standard way to solve for spectra, i.e., discretize and eigensolve. This is beautiful work by Boffi. (1 slide)
        \item However, even when you get everything right, you are approximating a dense set with points and you may not get it right. Disuss these issues and show an example (curl-curl). (1 slide)
        \item The solution is Matt's algorithm which works in infinite-dimensional setting. Explain this. (2 slides)
        \item Explain the research question: How can we solve the discrete problems in Matt's algorithm in an adaptive and multilevel way. Adaptive because we solve a subproblem for an approximate bound on the distance to the spectrum that may not be great. Multilevel because there are in theory infinite subproblems to solve and we want to do this efficiently (efficiently use, say, 1000 cores). (2 slides)
        \item Immediate next steps (what you are doing now and next week) and long term steps (what is your plan)? 

        \item Examiners will focus on asking motivation questions about the project structure in this setting. The hard hitting background questions will be in the viva. 

        \item Have a draft to Patrick by Friday. 
    \end{enumerate}


    \item The next meeting is June 2, at 12:30 p.m. UK time.

    \item I have a meeting with Matt Colbrook sometime on June 11.
    
\end{itemize}


\subsection{This Week's Plan}

\begin{enumerate}
    \item Answer the meeting questions I didnt answer well, i.e., why we consider different mesh structures in Boffi. 
    \item Curl-Curl Problem.
    \item Focus on the deliverables for the presentation. This includes picking examples, understanding the gist of Matt's algorithm, understanding the motivation and idea of using adaptivity and multilevel solves.
    \item Construct the slides before Friday!
\end{enumerate}


\section{Daily Logs}

\textbf{Monday} 

\begin{enumerate}
    \item Prepped and had a meeting with Patrick. Made the notes above. 
    \item Went back through Boffi and focused on the point of considering all the different meshes. It seems to be that the mesh introduces `bias' for certain eigenfunction combinations. This is less predictable for unstructured meshes. For structured meshes, computing convergence rates are cleaner, and the direction of the diagonal introduces a bias in the computed eigenfunctions (swap the diagonal direction and you reflect the eigenfunctions). To introduce more symmetry, we solve on criss-cross meshes. In doing so, the discrete eigenvalues approximating $\lambda = 5$ and $\lambda = 10$ coincide (up to machine precision). Furthermore, in this example, we see that on meshes that are too coarse, we can't expect the discrete eigenvalues to be in a one-to-one correspondence with the continuous ones. In solving on unstructured L-shaped meshes, we show that the convergence rate deteriorates for singular eigenfunctions. Furthermore, we explore the motivation for adaptivity near the re-entrant corner. Finally, we consider nonconforming elements and find that the eigenvalues are approximated from below rather than from above.
    
\end{enumerate}

\textbf{Tuesday}
\begin{enumerate}
    \item Today I focused on the presentation.
    \item I did a quick review and made outline slides.
    \item Then I filled in the introductory spectral slides and the Matt Colbrook slides. To create the Matt Colbrook slides, I spent a few hours picking out pieces of his work. I found it tricky to do this quickly and thoroughly because there is so much in his text and it is quite general.
\end{enumerate}

\textbf{Wednesday}
\begin{enumerate}
    \item Spent a little bit of time on the slides. I was up most of the night on an overnight flight (with the flu), so most of today was spent in bed.
\end{enumerate}

\textbf{Thursday}
\begin{enumerate}
    \item Finished up some slides.
    \item Read about and coded up the Maxwell problem.
    \item Edited the presentation.
\end{enumerate}

\textbf{Friday}
\begin{enumerate}
    \item Worked more on the Maxwell example. Created the criss-cross example for the presentation. Convinced myself that the eigenvalues are the eigenvalues of the Neumann Laplacian on $(0, \pi)^2.$
    \item Worked through the derivation of the spectra of the right and left shift operators on $\ell^2.$ This was way more involved than I had anticipated so I was glad to have worked through it.
    \item Added the operator example at the end of the presentation in case I get asked. 
    \item Converted my description of Matt's algorithm to a tikz diagram.
    \item Made some notes about what I need to remember about each slide when presenting. 
\end{enumerate}


\textbf{Sunday}
\begin{enumerate}
\item Worked on the presentation all day today. I really didn't think it would take me another full day. 
\item Started reading about mulitgrid methods incase I am asked about them in the presentation. 
\end{enumerate}

\textbf{Monday}
\begin{enumerate}
\item Actually read-through and learned about the Maxwell examlple in Boffi's text. Before I had only read what was needed for implementation.  
\item Read through important bits of Matt's Chapter 3. I did this as a less thorough read-through. There is alot to dive into here, I just wanted to get the general sense of the algorithm for the presentation. 
\item Worked on Patrick's suggested edits for the presentation. 
\item Started writing up the first finite element examples (Laplace and Maxwell). Hopefully I will finish this tomorrow. 
\end{enumerate}


\end{document}
